% !TEX program = xelatex
\documentclass[a4paper,10pt]{article}

\makeatletter
\let\@expl@@@mark@update@singlecol@structures@@\relax
\makeatother

\usepackage[a4paper, left=0.95cm, right=0.95cm, top=0.85cm, bottom=0.92cm, footskip=0.52cm]{geometry}
\usepackage{fontspec}
\usepackage{amsmath, amssymb}
\usepackage{longtable, array, colortbl}
\usepackage[table]{xcolor}
\usepackage{fancyhdr}
\usepackage{tikz}
\usetikzlibrary{calc, angles, quotes, arrows.meta, decorations.markings}

% Cross-platform font discovery
\makeatletter
\newcommand{\NGSFindFonts}{%
  \IfFileExists{../../../fonts/KhmerOSsiemreap.ttf}{%
    \def\NGSUserPath{../../../fonts/}\def\NGSSysPath{../../../fonts/}%
  }{%
    \IfFileExists{../../../../fonts/KhmerOSsiemreap.ttf}{%
      \def\NGSUserPath{../../../../fonts/}\def\NGSSysPath{../../../../fonts/}%
    }{%
      \IfFileExists{fonts/KhmerOSsiemreap.ttf}{%
        \def\NGSUserPath{fonts/}\def\NGSSysPath{fonts/}%
      }{%
        \def\NGSUserPath{C:/Users/user/AppData/Local/Microsoft/Windows/Fonts/}%
        \def\NGSSysPath{C:/Windows/Fonts/}%
      }%
    }%
  }%
}
\NGSFindFonts
\makeatother

\XeTeXlinebreaklocale "km"
\XeTeXlinebreakskip = 0pt plus 1pt minus 0.1pt

\setmainfont[
  Path = {\NGSUserPath},
  Extension = .ttf,
  UprightFont = KhmerOSsiemreap,
  BoldFont = KhmerOSsiemreap,
  BoldFeatures = {FakeBold=1.8},
  ItalicFont = KhmerOSsiemreap,
  ItalicFeatures = {FakeSlant=0.25}
]{KhmerOSsiemreap}

\newfontfamily\muol[Path={\NGSUserPath}]{KhmerOSmuollight.ttf}
\newfontfamily\Tacteing[Path={\NGSUserPath}]{TACTENG.TTF}
\newfontfamily\enfont[Path={\NGSSysPath}, Scale=1.0]{times.ttf}
\newfontfamily\enbold[Path={\NGSSysPath}, Scale=1.0]{timesbd.ttf}
\newfontfamily\enital[Path={\NGSSysPath}, Scale=1.0]{timesi.ttf}
\newfontfamily\calibrifont[Path={\NGSSysPath}, Scale=0.95]{calibri.ttf}
\newfontfamily\symfont[Path={\NGSSysPath}, Scale=0.95]{cour.ttf}

\DeclareTextFontCommand{\en}{\enfont}
\newcommand{\enbf}[1]{{\enbold #1}}
\newcommand{\enit}[1]{{\enital #1}}
\newcommand{\khmflourish}{{\Tacteing\fontsize{17pt}{17pt}\selectfont 3}}
\newcommand{\wcirc}{{\symfont\char"25CB}}

\definecolor{headgray}{HTML}{E6E6E6}
\definecolor{stepgray}{HTML}{D4D4D4}
\definecolor{fiveecolor}{HTML}{8B0000}

\pagestyle{fancy}
\fancyhf{}
\renewcommand{\headrulewidth}{0pt}
\renewcommand{\footrulewidth}{0pt}
\fancyfoot[R]{{\calibrifont\fontsize{9.5pt}{11pt}\selectfont\thepage}}

\setlength{\parindent}{0pt}
\setlength{\LTpre}{2pt}
\setlength{\LTpost}{2pt}
\setlength{\tabcolsep}{4.5pt}
\newcolumntype{L}[1]{>{\raggedright\arraybackslash}p{#1}}

\newcommand{\fiveE}[2]{{\color{fiveecolor}\enbf{#1} \enit{(#2)}}}
\newcommand{\timenote}[1]{{\color{fiveecolor}\textit{#1}}}
\newcommand{\stepheader}[1]{%
  \rowcolor{stepgray}%
  \multicolumn{3}{|c|}{\rule[-4pt]{0pt}{15pt}{\muol\fontsize{10pt}{14pt}\selectfont #1}} \\*%
  \hline
}
\newcommand{\cfig}[1]{\begin{center}\vspace{-2pt}#1\vspace{-3pt}\end{center}}

\newcommand{\tri}[1]{\triangle #1}
\newcommand{\unit}[1]{\,\text{\en{#1}}}
\newcommand{\cnum}[1]{\tikz[baseline=-0.58ex]\node[circle, fill=black, text=white, inner sep=0.85pt, minimum size=1.95ex, font=\enbold\scriptsize]{#1};}

\newcommand{\arc}[1]{\mbox{%
  \tikz[baseline=(arcX.base)]{%
    \node[inner sep=0pt, outer sep=0pt] (arcX) {$#1$};%
    \draw[black, line width=0.55pt] ($(arcX.north west)+(0.02,0.035)$) to[bend left=30] ($(arcX.north east)+(-0.02,0.035)$);%
  }%
}}

\newcommand{\tickone}[2]{%
  \coordinate (tkM) at ($(#1)!0.5!(#2)$);
  \draw[black, line width=0.7pt] ($(tkM)!0.11cm!90:(#2)$) -- ($(tkM)!0.11cm!-90:(#2)$);
}
\newcommand{\ticktwo}[2]{%
  \coordinate (tkM1) at ($(#1)!0.47!(#2)$);
  \coordinate (tkM2) at ($(#1)!0.53!(#2)$);
  \draw[black, line width=0.7pt] ($(tkM1)!0.11cm!90:(#2)$) -- ($(tkM1)!0.11cm!-90:(#2)$);
  \draw[black, line width=0.7pt] ($(tkM2)!0.11cm!90:(#2)$) -- ($(tkM2)!0.11cm!-90:(#2)$);
}

% ============================================================================
% ម៉ាក្រូរូបភាព TikZ ស-ខ្មៅ សមាមាត្រច្បាស់ (3.8 cm – 6.8 cm, អក្សរខ្មៅ 100%)
% ============================================================================

% ១. រូបបេសកកម្ម E1៖ មុខកាត់ទទឹងលូបេតុងមូល C(O, 5dm) និងផ្ទៃទឹក CD នៅចម្ងាយ OH = 3dm
\newcommand{\figEngageCIRCFour}{%
\begin{tikzpicture}[scale=0.50, every node/.style={font=\footnotesize, text=black}]
  \coordinate (O) at (0,0);
  \def\R{2.5}
  \def\yH{-1.5}
  \def\xC{2.0}
  \coordinate (A) at (0,\R);
  \coordinate (B) at (0,-\R);
  \coordinate (H) at (0,\yH);
  \coordinate (C) at (-\xC,\yH);
  \coordinate (D) at (\xC,\yH);
  % កម្រាស់បេតុងខាងក្រៅ និងផ្ទៃទឹកឆ្នូតស-ខ្មៅស្រាល
  \draw[black!55, line width=0.7pt] (O) circle (2.78);
  \begin{scope}
    \clip (O) circle (\R);
    \fill[black!10] (-2.6,-2.6) rectangle (2.6,\yH);
  \end{scope}
  \draw[black, line width=1.15pt] (O) circle (\R);
  \draw[black, line width=1.15pt] (C) -- (D);
  \draw[black, line width=0.9pt, densely dashed] (A) -- (B);
  \draw[black, line width=0.95pt] (O) -- (D);
  % សញ្ញាមុំកែងត្រង់ H
  \draw[black, line width=0.7pt] ($(H)+(0,0.26)$) -- ($(H)+(0.26,0.26)$) -- ($(H)+(0.26,0)$);
  \fill[black] (O) circle (2.2pt) (A) circle (2pt) (B) circle (2pt) (H) circle (2pt) (C) circle (2pt) (D) circle (2pt);
  \node[above left=0.5pt] at (O) {$O$};
  \node[above=1.5pt] at (A) {$A$};
  \node[below=1.5pt] at (B) {$B$};
  \node[below left=0.5pt] at (H) {$H$};
  \node[left=2pt] at (C) {$C$};
  \node[right=2pt] at (D) {$D$};
  \node[left=1pt] at (0,-0.72) {$3\unit{dm}$};
  \node[above right=-1pt] at (1.0,-0.72) {$5\unit{dm}$};
  \node at (0,-3.95) {ទទឹងផ្ទៃទឹក $CD = ?$};
\end{tikzpicture}}

% ២. រូប E2៖ ផ្ចិតឆ្លុះ និងអ័ក្សឆ្លុះ (ឧ.១-៣ ទំ.២០២) និងអង្កត់ផ្ចិតកែងនឹងអង្កត់ធ្នូ (ទំ.២០២-២០៣)
\newcommand{\figExploreCIRCFour}{%
\begin{tikzpicture}[scale=0.36, every node/.style={font=\scriptsize, text=black}]
  % ខាងឆ្វេង៖ ផ្ចិតឆ្លុះ O និងអ័ក្សឆ្លុះ d (ទំ.២០២)
  \begin{scope}[shift={(0,0)}]
    \coordinate (O) at (3.0,2.8);
    \def\R{2.35}
    \coordinate (M1) at ($(O)+(-\R,0)$);
    \coordinate (M2) at ($(O)+(\R,0)$);
    \coordinate (P) at ($(O)+(55:\R)$);
    \coordinate (Psym) at ($(O)+(-55:\R)$);
    \coordinate (Pcen) at ($(O)+(235:\R)$);
    \coordinate (K) at ($(P)!(O)!(Psym)$);
    \draw[black, line width=0.95pt] (O) circle (\R);
    \draw[black, line width=1.05pt] ($(M1)+(-0.35,0)$) -- ($(M2)+(0.35,0)$) node[right=0.5pt] {$d$};
    \draw[black, line width=0.85pt, densely dashed] (P) -- (Psym) (P) -- (Pcen);
    \draw[black, line width=0.65pt] ($(K)+(0,0.22)$) -- ($(K)+(-0.22,0.22)$) -- ($(K)+(-0.22,0)$);
    \tickone{O}{P}\tickone{O}{Pcen}
    \ticktwo{P}{K}\ticktwo{K}{Psym}
    \fill[black] (O) circle (2.2pt) (P) circle (2pt) (Psym) circle (2pt) (Pcen) circle (2pt) (K) circle (1.8pt);
    \node[above left=0.5pt] at (O) {$O$};
    \node[above right=0.5pt] at (P) {$B$};
    \node[below right=0.5pt] at (Psym) {$B_1$};
    \node[below left=0.5pt] at (Pcen) {$B'$};
    \node[align=center] at (3.0,-1.75) {ផ្ចិតឆ្លុះ $O$ ($OB=OB'$)\\និងអ័ក្សឆ្លុះ $d$ ($BK=B_1K$)};
  \end{scope}
  % ខាងស្ដាំ៖ អង្កត់ផ្ចិត AB ⊥ CD ត្រង់ H (គំរូទី១ ទំ.២០៣)
  \begin{scope}[shift={(9.8,0)}]
    \coordinate (O2) at (3.0,2.8);
    \def\R{2.35}
    \coordinate (A2) at ($(O2)+(0,\R)$);
    \coordinate (B2) at ($(O2)+(0,-\R)$);
    \coordinate (C2) at ($(O2)+(-145:\R)$);
    \coordinate (D2) at ($(O2)+(-35:\R)$);
    \coordinate (H2) at ($(C2)!(O2)!(D2)$);
    \draw[black, line width=0.95pt] (O2) circle (\R);
    \draw[black, line width=1.05pt] (A2)--(B2) (C2)--(D2);
    \draw[black, line width=0.8pt, densely dashed] (O2)--(C2) (O2)--(D2);
    \draw[black, line width=0.65pt] ($(H2)+(0,0.22)$) -- ($(H2)+(0.22,0.22)$) -- ($(H2)+(0.22,0)$);
    \tickone{O2}{C2}\tickone{O2}{D2}
    \ticktwo{C2}{H2}\ticktwo{H2}{D2}
    \fill[black] (O2) circle (2.2pt) (A2) circle (2pt) (B2) circle (2pt) (C2) circle (2pt) (D2) circle (2pt) (H2) circle (1.9pt);
    \node[right=1.5pt] at (O2) {$O$};
    \node[above=1pt] at (A2) {$A$};
    \node[below=1pt] at (B2) {$B$};
    \node[below left=0.5pt] at (C2) {$C$};
    \node[below right=0.5pt] at (D2) {$D$};
    \node[above left=0.5pt] at (H2) {$H$};
    \node[align=center] at (3.0,-1.75) {$AB \perp CD$ ត្រង់ $H$\\$\Rightarrow HC=HD,\; \arc{CB}=\arc{DB}$};
  \end{scope}
\end{tikzpicture}}

% ៣. រូប E4៖ គំរូទី២ (AB ∥ CD ⇒ ធ្នូ AC = BD ទំ.២០៣) និងលំហាត់ទី៤ (OQ = OR ⇒ AB = CD ទំ.២០៤)
\newcommand{\figModelCIRCFour}{%
\begin{tikzpicture}[scale=0.36, every node/.style={font=\scriptsize, text=black}]
  % គំរូទី២ ទំ.២០៣៖ AB ∥ CD កាត់ដោយអង្កត់ផ្ចិត MN ⊥ AB
  \begin{scope}[shift={(0,0)}]
    \coordinate (O) at (3.0,2.8);
    \def\R{2.35}
    \coordinate (M) at ($(O)+(0,\R)$);
    \coordinate (N) at ($(O)+(0,-\R)$);
    \coordinate (A) at ($(O)+(135:\R)$);
    \coordinate (B) at ($(O)+(45:\R)$);
    \coordinate (C) at ($(O)+(195:\R)$);
    \coordinate (D) at ($(O)+(-15:\R)$);
    \coordinate (H1) at ($(A)!(O)!(B)$);
    \coordinate (H2) at ($(C)!(O)!(D)$);
    \draw[black, line width=0.95pt] (O) circle (\R);
    \draw[black, line width=1.05pt] (A)--(B) (C)--(D);
    \draw[black, line width=0.85pt, densely dashed] (M)--(N);
    \draw[black, line width=1.5pt] (A) arc (135:195:\R);
    \draw[black, line width=1.5pt] (D) arc (-15:45:\R);
    \draw[black, line width=0.65pt] ($(H1)+(0,-0.20)$) -- ($(H1)+(0.20,-0.20)$) -- ($(H1)+(0.20,0)$);
    \draw[black, line width=0.65pt] ($(H2)+(0,0.20)$) -- ($(H2)+(0.20,0.20)$) -- ($(H2)+(0.20,0)$);
    \fill[black] (O) circle (2.1pt) (M) circle (2pt) (N) circle (2pt) (A) circle (2pt) (B) circle (2pt) (C) circle (2pt) (D) circle (2pt);
    \node[above right=1pt] at (O) {$O$};
    \node[above=1pt] at (M) {$M$};
    \node[below=1pt] at (N) {$N$};
    \node[above left=0.5pt] at (A) {$A$};
    \node[above right=0.5pt] at (B) {$B$};
    \node[left=1pt] at (C) {$C$};
    \node[right=1pt] at (D) {$D$};
    \node[align=center] at (3.0,-1.75) {គំរូទី២៖ $AB \parallel CD$\\$MN \perp AB \Rightarrow \arc{AC}=\arc{BD}$};
  \end{scope}
  % លំហាត់ទី៤ ទំ.២០៤៖ OQ ⊥ AB, OR ⊥ CD, OQ = OR
  \begin{scope}[shift={(10.2,0)}]
    \coordinate (O2) at (3.0,2.8);
    \def\R{2.35}
    \coordinate (A2) at ($(O2)+(165:\R)$);
    \coordinate (B2) at ($(O2)+(55:\R)$);
    \coordinate (C2) at ($(O2)+(-145:\R)$);
    \coordinate (D2) at ($(O2)+(-35:\R)$);
    \coordinate (Q2) at ($(A2)!(O2)!(B2)$);
    \coordinate (R2) at ($(C2)!(O2)!(D2)$);
    \draw[black, line width=0.95pt] (O2) circle (\R);
    \draw[black, line width=1.05pt] (A2)--(B2) (C2)--(D2);
    \draw[black, line width=0.85pt] (O2)--(Q2) (O2)--(R2);
    \draw[black, line width=0.65pt] ($(Q2)!0.20cm!(B2)$) -- ($(Q2)!0.20cm!(B2)!0.20cm!90:(B2)$) -- ($(Q2)!0.20cm!(O2)$);
    \draw[black, line width=0.65pt] ($(R2)+(0,0.20)$) -- ($(R2)+(0.20,0.20)$) -- ($(R2)+(0.20,0)$);
    \tickone{O2}{Q2}\tickone{O2}{R2}
    \fill[black] (O2) circle (2.2pt) (A2) circle (2pt) (B2) circle (2pt) (C2) circle (2pt) (D2) circle (2pt) (Q2) circle (1.9pt) (R2) circle (1.9pt);
    \node[right=1.5pt] at (O2) {$O$};
    \node[left=1pt] at (A2) {$A$};
    \node[above right=0.5pt] at (B2) {$B$};
    \node[below left=0.5pt] at (C2) {$C$};
    \node[below right=0.5pt] at (D2) {$D$};
    \node[below left=0.5pt] at (Q2) {$Q$};
    \node[above left=0.5pt] at (R2) {$R$};
    \node[align=center] at (3.0,-1.75) {លំហាត់ទី៤៖ $OQ=OR$\\$\Rightarrow x=-1,\; AB=CD=8$};
  \end{scope}
\end{tikzpicture}}

% ៤. រូប E5៖ ប្រតិបត្តិ ទំ.២០៣ (អង្កត់ផ្ចិត AD និង AB ∥ CD ⇒ AB = CD) និងលំហាត់ទី៨ ទំ.២០៤
\newcommand{\figPracticeCIRCFour}{%
\begin{tikzpicture}[scale=0.36, every node/.style={font=\scriptsize, text=black}]
  % ប្រតិបត្តិ ទំ.២០៣
  \begin{scope}[shift={(0,0)}]
    \coordinate (O) at (3.0,2.8);
    \def\R{2.35}
    \coordinate (A) at ($(O)+(180:\R)$);
    \coordinate (D) at ($(O)+(0:\R)$);
    \coordinate (B) at ($(O)+(110:\R)$);
    \coordinate (C) at ($(O)+(-70:\R)$);
    \draw[black, line width=0.95pt] (O) circle (\R);
    \draw[black, line width=1.05pt] (A)--(D) (A)--(B) (C)--(D);
    \draw[black, line width=0.8pt, densely dashed] (O)--(B) (O)--(C);
    \tickone{O}{A}\tickone{O}{D}\tickone{O}{B}\tickone{O}{C}
    \fill[black] (O) circle (2.2pt) (A) circle (2pt) (B) circle (2pt) (C) circle (2pt) (D) circle (2pt);
    \node[above right=0.5pt] at (O) {$O$};
    \node[left=1pt] at (A) {$A$};
    \node[right=1pt] at (D) {$D$};
    \node[above=1pt] at (B) {$B$};
    \node[below=1pt] at (C) {$C$};
    \node[align=center] at (3.0,-1.75) {ប្រតិបត្តិ៖ អង្កត់ផ្ចិត $AD$\\$AB \parallel CD \Rightarrow AB=CD$};
  \end{scope}
  % លំហាត់ទី៨ ទំ.២០៤ (កន្លះរង្វង់ អង្កត់ផ្ចិត AB, AP ⊥ CD, BS ⊥ CD ⇒ PC = DS)
  \begin{scope}[shift={(9.8,0)}]
    \coordinate (O2) at (3.0,1.2);
    \def\R{2.35}
    \coordinate (A2) at ($(O2)+(-\R,0)$);
    \coordinate (B2) at ($(O2)+(\R,0)$);
    \coordinate (C2) at ($(O2)+(130:\R)$);
    \coordinate (D2) at ($(O2)+(40:\R)$);
    \coordinate (P2) at ($(C2)!(A2)!(D2)$);
    \coordinate (S2) at ($(C2)!(B2)!(D2)$);
    \coordinate (H2) at ($(C2)!(O2)!(D2)$);
    \draw[black, line width=0.95pt] (B2) arc (0:180:\R) -- cycle;
    \draw[black, line width=1.0pt] ($(P2)!-0.25cm!(S2)$) -- ($(S2)!-0.25cm!(P2)$);
    \draw[black, line width=0.85pt] (A2)--(P2) (B2)--(S2);
    \draw[black, line width=0.8pt, densely dashed] (O2)--(H2);
    \fill[black] (O2) circle (2.1pt) (A2) circle (2pt) (B2) circle (2pt) (C2) circle (2pt) (D2) circle (2pt) (P2) circle (1.8pt) (S2) circle (1.8pt) (H2) circle (1.8pt);
    \node[below=1.5pt] at (O2) {$O$};
    \node[below left=0.5pt] at (A2) {$A$};
    \node[below right=0.5pt] at (B2) {$B$};
    \node[above=1.5pt] at (C2) {$C$};
    \node[above=1.5pt] at (D2) {$D$};
    \node[above left=0.5pt] at (P2) {$P$};
    \node[above right=0.5pt] at (S2) {$S$};
    \node[below right=0pt] at (H2) {$H$};
    \node[align=center] at (3.0,-1.75) {លំហាត់ទី៨៖ $OH \perp CD$\\$\Rightarrow HP=HS,\; HC=HD$};
  \end{scope}
\end{tikzpicture}}

\begin{document}
\fontsize{10pt}{15.0pt}\selectfont

% ==================== ក្បាលរដ្ឋបាលផ្លូវការ ====================
{\centering
  {\muol\fontsize{11pt}{16pt}\selectfont ព្រះរាជាណាចក្រកម្ពុជា}\\[2pt]
  {\muol\fontsize{11pt}{16pt}\selectfont ជាតិ សាសនា ព្រះមហាក្សត្រ}\\[-2pt]
  \khmflourish\par
}
\vspace{-4pt}
{\muol\fontsize{10pt}{15.2pt}\selectfont
មន្ទីរអប់រំ យុវជន និងកីឡា ខេត្តកំពង់ធំ\\
ការិយាល័យអប់រំ យុវជន និងកីឡា ស្រុកកំពង់ស្វាយ\\
អនុវិទ្យាល័យព្រះក្រសាំង\par}

\vspace{2pt}
{\centering{\muol\fontsize{12pt}{17pt}\selectfont កិច្ចតែងការបង្រៀន}\par}
\vspace{2pt}

% ==================== ព័ត៌មានទូទៅ ====================
\hspace*{3.3cm}
\begin{tabular}{@{}l@{\hspace{6pt}}l@{\hspace{6pt}}l@{}}
  - & មុខវិជ្ជា & ៖ គណិតវិទ្យា ថ្នាក់ទី៨ \\
  - & ចំណងជើងមេរៀន & ៖ មេរៀនទី១៥ រង្វង់ \\
  - & ចំណងជើងរង & ៖ ២.២ លក្ខណៈឆ្លុះនៃរង្វង់ និងលំហាត់សំយោគ (ទំព័រ ២០២--២០៤) \\
  - & រយៈពេល & ៖ ១ ម៉ោងសិក្សា (៥០ នាទី) \\
  - & បង្រៀនដោយ & ៖ ហុក ប៊ុនថែន \\
  - & កាលបរិច្ឆេទ & ៖ ថ្ងៃអង្គារ ទី២២ ខែធ្នូ ឆ្នាំ២០២៦
\end{tabular}

\vspace{4pt}
% ==================== I. វត្ថុបំណង ====================
{\muol\fontsize{10pt}{15pt}\selectfont \enbf{I.} វត្ថុបំណង}\par\vspace{1pt}
\begin{tabular}{@{}p{0.45cm}@{}p{2.35cm}@{}p{0.35cm}@{}p{15.65cm}@{}}
  \wcirc & វិជ្ជាសម្បទា & ៖ & សិស្សប្រាប់បាននូវផ្ចិតឆ្លុះ អ័ក្សឆ្លុះនៃរង្វង់ ទ្រឹស្ដីបទអង្កត់ផ្ចិតកែងនឹងអង្កត់ធ្នូ អង្កត់ធ្នូស្ថិតនៅចម្ងាយស្មើគ្នាពីផ្ចិត និងអង្កត់ធ្នូពីរស្របគ្នាក្នុងរង្វង់បានត្រឹមត្រូវ ។ \\[2pt]
  \wcirc & បំណិនសម្បទា & ៖ & សិស្សប្រើលក្ខណៈឆ្លុះនៃរង្វង់ និងធនាគារហេតុផល \cnum{1}--\cnum{5} ស្រាយបញ្ជាក់លក្ខណៈអង្កត់ធ្នូ--ធ្នូ និងគណនាប្រវែងអង្កត់ធ្នូតាមទ្រឹស្ដីបទពីតាករ និងសមីការដឺក្រេទី១ លើសន្លឹកកិច្ចការបានត្រឹមត្រូវ ។ \\[2pt]
  \wcirc & ចរិយាសម្បទា & ៖ & សិស្សមានទម្លាប់ពិភាក្សាផ្ទៀងផ្ទាត់អំណះអំណាងជាដៃគូ (\en{Think-Pair-Share}) ហ្មត់ចត់លើលក្ខខណ្ឌកាត់តាមផ្ចិត និងឱ្យតម្លៃការអនុវត្តធរណីមាត្ររង្វង់ក្នុងវិស្វកម្មសំណង់ ។
\end{tabular}

\vspace{3pt}
% ==================== II. សម្ភារៈឧបទេស ====================
{\muol\fontsize{10pt}{15pt}\selectfont \enbf{II.} សម្ភារៈឧបទេស}\par\vspace{1pt}
\hspace*{0.35cm}$\bullet$ \textbf{សម្រាប់គ្រូ៖} សៀវភៅសិក្សាគោលគណិតវិទ្យាថ្នាក់ទី៨ (ទំ.២០២--២០៤) ក្ដារខៀន ដីស ដែកឈាន និងបន្ទាត់កែងធំ (គ្មានស្លាយ \en{LCD}) ។\par
\hspace*{0.35cm}$\bullet$ \textbf{សម្រាប់សិស្ស៖} សន្លឹកកិច្ចការសិស្ស ស-ខ្មៅ (\en{Worksheet\_CIRC4} ២ ទំព័រពេញ) សៀវភៅគោល បន្ទាត់កែង និងដែកឈាន ។

\vspace{3pt}
% ==================== III. វិធីសាស្រ្តបង្រៀន ====================
{\muol\fontsize{10pt}{15pt}\selectfont \enbf{III.} វិធីសាស្រ្តបង្រៀន}\par\vspace{1pt}
\begin{tabular}{@{}p{0.45cm}@{}p{2.55cm}@{}p{0.35cm}@{}p{15.45cm}@{}}
  $\bullet$ & គោលវិធី & ៖ & សិស្សមជ្ឈមណ្ឌល (បង្រៀនដោយមិនប្រើស្លាយ \en{LCD} យកសន្លឹកកិច្ចការរុករក និងសំយោគជាគោល) \\
  $\bullet$ & យុទ្ធវិធីបង្រៀន & ៖ & គំរូ \en{5Es} រួមផ្សំការរុករកតាមអ័ក្សឆ្លុះ ការពិភាក្សាដៃគូ (\en{Think-Pair-Share}) និងធនាគារហេតុផល \cnum{1}--\cnum{5}
\end{tabular}

\vspace{3pt}
% ==================== IV. ដំណើរការបង្រៀន ====================
{\muol\fontsize{10pt}{15pt}\selectfont \enbf{IV.} ដំណើរការបង្រៀន}\par\vspace{2pt}

\begin{longtable}{|L{5.50cm}|L{7.90cm}|L{4.80cm}|}
\hline
\rowcolor{headgray}
\multicolumn{1}{|c|}{\rule[-4pt]{0pt}{15pt}{\muol សកម្មភាពគ្រូ}} &
\multicolumn{1}{c|}{{\muol ខ្លឹមសារមេរៀន}} &
\multicolumn{1}{c|}{{\muol សកម្មភាពសិស្ស}} \\
\hline
\endhead

% ==================== ជំហានទី១ (2min) ====================
\stepheader{ជំហានទី១ ៖ រដ្ឋបាលថ្នាក់ (២ នាទី)}
$\bullet$ ពិនិត្យអវត្តមាន សណ្ដាប់ធ្នាប់ និងការត្រៀមបន្ទាត់កែង ដែកឈានរបស់សិស្ស\par
$\bullet$ ចែកសន្លឹកកិច្ចការ \en{Worksheet\_CIRC4} (១ សន្លឹកសងខាង) ដល់សិស្សគ្រប់គ្នា
&
$\bullet$ ពិនិត្យវត្តមាន វិន័យ និងសម្ភារៈសិក្សា (សន្លឹកកិច្ចការ បន្ទាត់កែង និងដែកឈាន)
&
$\circ$ ប្រធានថ្នាក់រាយការណ៍អវត្តមាន\par
$\circ$ ទទួលសន្លឹកកិច្ចការ និងសរសេរឈ្មោះ \\[3pt]
\hline

% ==================== ជំហានទី២ (9min: កែ ≤3 + រំលឹក 3 + E1 3) ====================
\stepheader{ជំហានទី២ ៖ កែកិច្ចការផ្ទះ រំលឹកមេរៀនចាស់ និងភ្ជាប់មេរៀនថ្មី (៩ នាទី)}
\timenote{(កែកិច្ចការផ្ទះ $\le 3\unit{min}$)}\par\vspace{1pt}
$\bullet$ ហៅសិស្ស ២ នាក់ឡើងបង្ហាញសម្រាយបញ្ជាក់សង្ខេបនៃលំហាត់ទី២ និងទី៥ (ទំ.២០៤)
&
\textbf{កំណែកិច្ចការផ្ទះ (ទំ.២០៤)៖}\par\vspace{1pt}
១. \textbf{លំហាត់ទី២៖} $\arc{AD} = \arc{CB} \Rightarrow \mathbf{AD = CB}$ (ទ្រឹស្ដីបទ ២) ។ បើ $\arc{AC} = \arc{BD}$ បូកធ្នូរួម $\arc{CD}$ បាន $\arc{AD} = \arc{BC} \Rightarrow AD = BC$ ។\par\vspace{1pt}
២. \textbf{លំហាត់ទី៥៖} $AB = CD \Rightarrow \arc{AB} = \arc{CD}$ ; បូកធ្នូរួម $\arc{BC} \Rightarrow \mathbf{\arc{AC} = \arc{BD}}$ ។
&
$\circ$ ឡើងកែលើក្ដារខៀន និងផ្ទៀងផ្ទាត់ចម្លើយក្នុងសៀវភៅ \\[3pt]

\timenote{(រំលឹកមេរៀនទាក់ទង $3\unit{min}$)}\par\vspace{1pt}
$\bullet$ សួររំលឹកនិយមន័យផ្ចិតឆ្លុះ អ័ក្សឆ្លុះ និងទ្រឹស្ដីបទពីតាករក្នុងត្រីកោណកែង
&
\textbf{រំលឹកបម្លែងឆ្លុះ និងត្រីកោណកែង៖}\par\vspace{1pt}
$\bullet$ \textbf{អ័ក្សឆ្លុះ $d$៖} បន្ទាត់ $d$ ជាមេដ្យាទ័រនៃអង្កត់ភ្ជាប់ចំណុចឆ្លុះ ($d \perp BB_1$ ត្រង់កណ្ដាល $K$) ។
&
$\circ$ ឆ្លើយរំលឹកលក្ខណៈផ្ចិតឆ្លុះ និងអ័ក្សឆ្លុះ \\[3pt]

$\bullet$ សួរបន្តអំពីករណីប៉ុនគ្នានៃត្រីកោណកែង និងទ្រឹស្ដីបទពីតាករ
&
$\bullet$ \textbf{ត្រីកោណកែងប៉ុនគ្នា៖} បើអ៊ីប៉ូតេនុស និងជ្រុងកែងមួយស្មើគ្នារៀងគ្នា នោះត្រីកោណកែងទាំងពីរប៉ុនគ្នា ។\par\vspace{1pt}
$\bullet$ \textbf{ពីតាករក្នុង $\tri{OHC}$ កែងត្រង់ $H$៖} $OH^2 + HC^2 = OC^2$ ។
&
$\circ$ រំលឹកករណីប៉ុនគ្នានៃត្រីកោណកែង និងរូបមន្តពីតាករ \\[3pt]

\fiveE{E1 - Engage}{3 min}\par\vspace{1pt}
$\bullet$ ឱ្យសិស្សពិនិត្យ \textbf{ផ្នែកទី១} លើសន្លឹកកិច្ចការ៖ មុខកាត់ទទឹងនៃ \textbf{លូបេតុងមូល} កាំក្នុង $r = 5\unit{dm}$ មានផ្ទៃទឹកដេក $CD$ ស្ថិតនៅចម្ងាយកែង $OH = 3\unit{dm}$ ក្រោមផ្ចិត $O$
&
\textbf{បេសកកម្ម E1៖ ទទឹងផ្ទៃទឹកក្នុងលូបេតុងមូល ($r = 5\unit{dm}$)}\par\vspace{1pt}
$\bullet$ សំណួរបំផុស៖ (១) តើអង្កត់ផ្ចិតឈរ $AB \perp CD$ ចែកទទឹងផ្ទៃទឹក $CD$ ជាពីរផ្នែកស្មើគ្នា ($HC = HD$) ដែរឬទេ? (២) តើទទឹងផ្ទៃទឹកពេញ $CD$ ស្មើ $4\unit{dm}$ ឬ $8\unit{dm}$? (ទុកផ្ទៀងនៅ E5)
&
$\circ$ សង្កេតរូបលូបេតុងមូល និងគូសធីកទស្សន៍ទាយទាំង ២ សំណួរលើសន្លឹកកិច្ចការ \\[3pt]

&
\cfig{\figEngageCIRCFour}
& \\[3pt]
\hline

% ==================== ជំហានទី៣ (17min: E2 11 + E3 6) ====================
\stepheader{ជំហានទី៣ ៖ រុករកលក្ខណៈឆ្លុះនៃរង្វង់ និងទ្រឹស្ដីបទអង្កត់ផ្ចិតកែងនឹងអង្កត់ធ្នូ (១៧ នាទី)}
\fiveE{E2 - Explore}{11 min}\par\vspace{1pt}
$\bullet$ ណែនាំដៃគូសិស្សប្រើបន្ទាត់កែងពិនិត្យរូបទី១ ក្នុង \textbf{ផ្នែកទី២} (ឧទាហរណ៍ទី១--៣ ទំ.២០២)៖ ប្រៀបធៀប $OB$ នឹង $OB'$ និងចំណុចឆ្លុះ $B_1$ ធៀបនឹងអង្កត់ផ្ចិត $d$
&
\textbf{សកម្មភាពទី១៖ ផ្ចិតឆ្លុះ និងអ័ក្សឆ្លុះនៃរង្វង់ (ទំ.២០២)}\par\vspace{1pt}
-- គ្រប់អង្កត់ផ្ចិត $BB'$ មាន $OB = OB' = r$ $\Rightarrow$ ផ្ចិត $O$ ជា \textbf{ផ្ចិតឆ្លុះ} នៃរង្វង់ ។\par\vspace{1pt}
-- បន្ទាត់ $d$ កាត់តាមផ្ចិត $O$ ជាមេដ្យាទ័រនៃអង្កត់ $BB_1$ ($BK = B_1K$) $\Rightarrow$ \textbf{អង្កត់ផ្ចិតគ្រប់ខ្សែជាអ័ក្សឆ្លុះ} នៃរង្វង់ ។
&
$\circ$ ប្រើបន្ទាត់កែងផ្ទៀងផ្ទាត់ផ្ចិតឆ្លុះ $O$ និងអ័ក្សឆ្លុះលើសន្លឹកកិច្ចការ \\[3pt]

$\bullet$ ឱ្យដៃគូពិនិត្យរូបទី២ (អង្កត់ផ្ចិត $AB \perp CD$ ត្រង់ $H$) និងពិភាក្សា \textbf{អន្ទាក់លក្ខខណ្ឌកាត់ផ្ចិត} (បើអង្កត់ធ្នូធម្មតាចែកអង្កត់ធ្នូមួយទៀតជាពីរស្មើគ្នា តើវាកែងជានិច្ចឬទេ?)
&
\textbf{សកម្មភាពទី២៖ អង្កត់ផ្ចិតកែងនឹងអង្កត់ធ្នូ និងអន្ទាក់មិនកាត់ផ្ចិត}\par\vspace{1pt}
-- ក្នុង $\tri{OHC}$ និង $\tri{OHD}$ កែងត្រង់ $H$ មាន $OC=OD=r$ និង $OH$ រួម $\Rightarrow \tri{OHC} \cong \tri{OHD} \Rightarrow \mathbf{HC = HD}$ និង $\angle BOC = \angle BOD \Rightarrow \mathbf{\arc{CB} = \arc{DB}}$ ។\par\vspace{1pt}
-- \textbf{អន្ទាក់គោលគំនិត៖} បើអង្កត់មួយ \textbf{មិនកាត់តាមផ្ចិត $O$} ទោះវាចែក $CD$ ជាពីរស្មើគ្នាក៏ដោយ ក៏វា \textbf{មិនកែង} នឹង $CD$ ដែរ!
&
$\circ$ រកឃើញថាត្រូវតែជាអង្កត់ផ្ចិត (កាត់តាម $O$) ទើបចែកអង្កត់ធ្នូមិនកាត់ផ្ចិតជាពីរស្មើគ្នាដោយកែង \\[3pt]

&
\cfig{\figExploreCIRCFour}
& \\[3pt]

\fiveE{E3 - Explain}{6 min}\par\vspace{1pt}
$\bullet$ សំយោគលក្ខណៈផ្ចិតឆ្លុះ និងអ័ក្សឆ្លុះនៃរង្វង់តាមសៀវភៅគោល (ទំ.២០២)
&
\textbf{សំយោគលក្ខណៈឆ្លុះ និងទ្រឹស្ដីបទគន្លឹះ (ទំ.២០២--២០៣)៖}\par\vspace{1pt}
$\bullet$ \textbf{លក្ខណៈឆ្លុះ៖} ផ្ចិត $O$ ជាផ្ចិតឆ្លុះ ហើយអង្កត់ផ្ចិតគ្រប់ខ្សែជាអ័ក្សឆ្លុះនៃរង្វង់ ។
&
$\circ$ បំពេញប្រអប់លក្ខណៈឆ្លុះក្នុង \textbf{ផ្នែកទី៣} នៃសន្លឹកកិច្ចការ \\[3pt]

$\bullet$ សំយោគទ្រឹស្ដីបទទី១ ស្ដីពីអង្កត់ផ្ចិតកែងនឹងអង្កត់ធ្នូ និងលក្ខណៈច្រាស (ទំ.២០២--២០៣)
&
$\bullet$ \textbf{ទ្រឹស្ដីបទ ១៖} អង្កត់ផ្ចិតកែងនឹងអង្កត់ធ្នូមួយ ចែកអង្កត់ធ្នូ និងធ្នូដែលវាកាត់ជាពីរផ្នែកស្មើគ្នា ($AB \perp CD \Rightarrow HC=HD,\; \arc{CB}=\arc{DB}$) និងច្រាសមកវិញ ។
&
$\circ$ បំពេញទ្រឹស្ដីបទទី១ និងកត់សម្គាល់លក្ខខណ្ឌកាត់តាមផ្ចិត $O$ \\[3pt]

$\bullet$ បន្តសំយោគទ្រឹស្ដីបទចម្ងាយពីផ្ចិត អង្កត់ធ្នូស្របគ្នា និងបង្ហាញ \textbf{ធនាគារហេតុផល \cnum{1}--\cnum{5}}
&
$\bullet$ \textbf{ទ្រឹស្ដីបទ ២៖} អង្កត់ធ្នូពីរស្មើគ្នា លុះត្រាតែស្ថិតនៅចម្ងាយកែងស្មើគ្នាពីផ្ចិត ($OQ \perp AB,\; OR \perp CD:\; OQ = OR \Leftrightarrow AB = CD$) ។\par\vspace{1pt}
$\bullet$ \textbf{ទ្រឹស្ដីបទ ៣៖} អង្កត់ធ្នូពីរស្របគ្នាក្នុងរង្វង់ កាត់ធ្នូចន្លោះប៉ុនគ្នា ($AB \parallel CD \Rightarrow \arc{AC} = \arc{BD}$) ។\par\vspace{1pt}
$\bullet$ \textbf{ធនាគារហេតុផល៖} \cnum{1} កាំរង្វង់ស្មើគ្នា $\;\vert\;$ \cnum{2} ត្រីកោណកែងប៉ុនគ្នា $\;\vert\;$ \cnum{3} អង្កត់ផ្ចិត $\perp$ អង្កត់ធ្នូ $\;\vert\;$ \cnum{4} ចម្ងាយកែងស្មើគ្នា $\Leftrightarrow$ អង្កត់ធ្នូស្មើគ្នា $\;\vert\;$ \cnum{5} ទ្រឹស្ដីបទពីតាករ ។
&
$\circ$ កត់ត្រាទ្រឹស្ដីបទទាំងអស់ និងធនាគារហេតុផល \cnum{1}--\cnum{5} \\[3pt]

$\bullet$ ឱ្យសិស្សសង្ខេបជាតារាង \textbf{សម្មតិកម្ម--សន្និដ្ឋាន} និងឆ្លើយពិត/មិនពិតលើសន្លឹកកិច្ចការ
&
{\renewcommand{\arraystretch}{1.12}\small
\begin{tabular}{|p{3.75cm}|p{3.55cm}|}
\hline
\cellcolor{headgray}\centering\textbf{សម្មតិកម្មក្នុង $C(O,r)$} &
\cellcolor{headgray}\centering\arraybackslash\textbf{សន្និដ្ឋាន / វិបាក} \\[2pt]
\hline
អង្កត់ផ្ចិត $AB \perp CD$ ត្រង់ $H$ & $HC = HD,\; \arc{CB} = \arc{DB}$ (\cnum{3}) \\[1.5pt]
\hline
$OQ \perp AB,\; OR \perp CD,\; OQ=OR$ & $AB = CD$ និងច្រាសមកវិញ (\cnum{4}) \\[1.5pt]
\hline
អង្កត់ធ្នូពីរស្របគ្នា $AB \parallel CD$ & ធ្នូចន្លោះ $\arc{AC} = \arc{BD}$ \\[1.5pt]
\hline
\end{tabular}}
&
$\circ$ ផ្ទៀងផ្ទាត់តារាងសម្មតិកម្ម--សន្និដ្ឋាន និងឆ្លើយពិត/មិនពិត \\[3pt]
\hline

% ==================== ជំហានទី៤ (20min: E4 10 + E5 10) ====================
\stepheader{ជំហានទី៤ ៖ ពង្រឹងចំណេះដឹងតាមសម្រាយបញ្ជាក់ សំយោគពីជគណិត និងវាយតម្លៃ (២០ នាទី)}
\fiveE{E4 - Elaborate}{10 min}\par\vspace{1pt}
$\bullet$ ដឹកនាំសិស្សបំពេញសម្រាយបញ្ជាក់ \textbf{លំហាត់គំរូទី១} និង \textbf{លំហាត់គំរូទី២ (ទំ.២០៣)} លើ \textbf{ផ្នែកទី៤} ដោយប្រើធនាគារហេតុផល
&
\textbf{១. លំហាត់គំរូទី១ (ទំ.២០៣ — $AB \perp CD \Rightarrow HC=HD,\; \arc{CB}=\arc{DB}$)៖}\par
-- ក្នុង $\tri{OHC}$ និង $\tri{OHD}$ កែងត្រង់ $H$ មាន $OC=OD=r$ (\cnum{1}) និង $OH$ រួម $\Rightarrow \tri{OHC} \cong \tri{OHD}$ (\cnum{2}) $\Rightarrow \mathbf{HC=HD}$ និង $\angle BOC = \angle BOD \Rightarrow \mathbf{\arc{CB}=\arc{DB}}$ ។\par\vspace{1pt}
\textbf{២. លំហាត់គំរូទី២ (ទំ.២០៣ — $AB \parallel CD \Rightarrow \arc{AC} = \arc{BD}$)៖}\par
-- គូសអង្កត់ផ្ចិត $MN \perp AB \Rightarrow MN \perp CD$ ដែរ ។ តាម \cnum{3} បាន $\arc{MA} = \arc{MB}$ និង $\arc{MC} = \arc{MD}$ $\Rightarrow \arc{MC} - \arc{MA} = \arc{MD} - \arc{MB} \Rightarrow \mathbf{\arc{AC} = \arc{BD}}$ ។
&
$\circ$ គូសអង្កត់ផ្ចិតជំនួយ $MN \perp AB$ និងបំពេញជំហានសម្រាយបញ្ជាក់ \\[3pt]

$\bullet$ ឱ្យដៃគូដោះស្រាយ \textbf{លំហាត់ទី៤ (ទំ.២០៤ — សំយោគពីជគណិត)} ដោយប្រើលក្ខណៈចម្ងាយស្មើគ្នាពីផ្ចិត \cnum{4}
&
\textbf{៣. លំហាត់ទី៤ (ទំ.២០៤ — $OQ \perp AB,\; OR \perp CD,\; OQ = OR$)៖}\par
-- គេឱ្យ $AB = 6x + 14$ និង $CD = -4x + 4$ ។ ដោយ $OQ = OR$ តាម \cnum{4} បាន $AB = CD \Rightarrow 6x + 14 = -4x + 4 \Rightarrow 10x = -10 \Rightarrow \mathbf{x = -1}$ ។\par
-- ជំនួស $x = -1$៖ $AB = 6(-1) + 14 = \mathbf{8}$ និង $CD = -4(-1) + 4 = \mathbf{8}$ ។
&
$\circ$ បង្កើតសមីការ $6x+14 = -4x+4$ រក $x=-1$ និងផ្ទៀងផ្ទាត់ $AB=CD=8$ \\[3pt]

&
\cfig{\figModelCIRCFour}
& \\[3pt]

\fiveE{E5 - Evaluate}{10 min}\par\vspace{1pt}
$\bullet$ ប្រើទម្រង់ \en{Think-Pair-Share} ឱ្យសិស្សគិតម្នាក់ឯង ១ នាទី លើចំណោទ \textbf{ប្រតិបត្តិ (ទំ.២០៣)} ក្នុង \textbf{ផ្នែកទី៥}
&
\textbf{១. ប្រតិបត្តិ (ទំ.២០៣)៖} ក្នុងរង្វង់ផ្ចិត $O$ មានអង្កត់ផ្ចិត $AD$ និងអង្កត់ធ្នូ $AB \parallel CD$ ។ ចូរស្រាយបញ្ជាក់ថា $AB = CD$ ៖
&
$\circ$ គិតម្នាក់ឯងស្វែងរកគន្លឹះភ្ជាប់អង្កត់ផ្ចិត $AD$ និង $AB \parallel CD$ \\[3pt]

$\bullet$ ឱ្យដៃគូពិភាក្សាផ្ទៀងផ្ទាត់ និងឡើងបង្ហាញសម្រាយបញ្ជាក់លើក្ដារខៀន
&
-- ដោយ $AB \parallel CD \Rightarrow \arc{AC} = \arc{BD}$ (តាមគំរូទី២) ។ តែ $AD$ ជាអង្កត់ផ្ចិត ($\arc{AB}+\arc{BD} = 180^\circ = \arc{AC}+\arc{CD}$) $\Rightarrow \arc{AB} = 180^\circ - \arc{BD} = 180^\circ - \arc{AC} = \arc{CD} \Rightarrow \mathbf{AB = CD}$! (ឬតាម $\tri{OAB} \cong \tri{ODC}$) ។
&
$\circ$ ពិភាក្សាដៃគូ និងសរសេរសម្រាយបញ្ជាក់លើសន្លឹកកិច្ចការ \\[3pt]

&
\cfig{\figPracticeCIRCFour}
& \\[3pt]

$\bullet$ ពិនិត្យចម្លើយបិទបេសកកម្មលូបេតុងមូល E1 ជាមួយសិស្សទូទាំងថ្នាក់
&
$\bullet$ \textbf{បិទបេសកកម្ម E1៖} ក្នុង $\tri{OHC}$ កែងត្រង់ $H$ តាម \cnum{5} បាន $HC = \sqrt{5^2 - 3^2} = \sqrt{16} = 4\unit{dm}$ ។ តាម \cnum{3} ទទឹងផ្ទៃទឹកពេញគឺ $CD = 2 \times HC = \mathbf{8\unit{dm}}$!
&
$\circ$ គណនាពីតាករ និងផ្ទៀងផ្ទាត់ថាមិនភ្លេចគុណនឹង $2$ ($CD = 8\unit{dm}$) \\[3pt]

$\bullet$ ធ្វើតេស្តរហ័ស ៦០ វិនាទី (\en{Exit Check}) ដោយឱ្យសិស្សលើកម្រាមដៃ ១--៤
&
$\bullet$ \textbf{Exit Check៖} ក្នុង $C(O, 10\unit{cm})$ អង្កត់ធ្នូ $AB$ ស្ថិតនៅចម្ងាយកែង $OH = 6\unit{cm}$ ពីផ្ចិត $O$ ។ ប្រវែងអង្កត់ធ្នូ $AB$ គឺ៖ A. $8\unit{cm}$ $\;\vert\;$ \textbf{B. $16\unit{cm}$ (ត្រូវ)} $\;\vert\;$ C. $12\unit{cm}$ $\;\vert\;$ D. $20\unit{cm}$ ។
&
$\circ$ លើកម្រាមដៃទី២ ជ្រើសចម្លើយត្រឹមត្រូវ \textbf{B} ($2\sqrt{10^2-6^2}=16\unit{cm}$) \\[3pt]
\hline

% ==================== ជំហានទី៥ (≤2min) ====================
\stepheader{ជំហានទី៥ ៖ កិច្ចការផ្ទះសំយោគបញ្ចប់មេរៀនទី១៥ និងបណ្ដាំផ្ញើ ($\le$ ២ នាទី)}
\timenote{(ចំណាយពេល $\le 2\unit{min}$)}\par\vspace{1pt}
$\bullet$ ដាក់កិច្ចការផ្ទះសំយោគបញ្ចប់មេរៀន៖ \textbf{លំហាត់ទី៨, ទី៩ និងទី១០ (ទំ.២០៤)}\par
$\bullet$ ផ្ដាំផ្ញើឱ្យសិស្សរំលឹកមេរៀនទី១៤ និងទី១៥ ទាំងមូលដើម្បីត្រៀមប្រឡងបញ្ចប់ជំពូក
&
\textbf{កិច្ចការផ្ទះសំយោគ (សៀវភៅសិក្សាគោល ទំព័រ ២០៤)៖}\par
១. \textbf{លំហាត់ទី ៨៖} កន្លះរង្វង់អង្កត់ផ្ចិត $AB$ អង្កត់ធ្នូ $CD$ មាន $AP \perp CD$ និង $BS \perp CD$ ។ បង្ហាញថា $PC = DS$ (ណែនាំ៖ គូស $OH \perp CD \Rightarrow HC=HD$ និង $HP=HS$) ។\par
២. \textbf{លំហាត់ទី ៩៖} ចំណុច $P$ ក្នុង $C(O,r)$ ។ បង្ហាញថាអង្កត់ធ្នូកាត់តាម $P$ ហើយកែងនឹង $OP$ មានប្រវែងខ្លីជាងគេ ។\par
៣. \textbf{លំហាត់ទី ១០ (ពង្រីក)៖} អង្កត់ធ្នូ $AB > CD$ កាត់គ្នាត្រង់ $K$ ក្រៅរង្វង់ ($H, P$ កណ្ដាល $AB, CD$ និង $M$ កណ្ដាល $OK$) ៖ បង្ហាញថា $MO=MK=MH=MP$ និង $KH > KP$ ។
&
$\circ$ កត់ត្រាកិច្ចការផ្ទះសំយោគទាំង ៣ លំហាត់\par
$\circ$ ទទួលបណ្ដាំផ្ញើ និងជម្រាបលាគ្រូ \\[3pt]
\hline
\end{longtable}

\vspace{6pt}
% ==================== ហត្ថលេខា ====================
\hfill
\begin{minipage}{8.6cm}
  \centering
  ថ្ងៃអង្គារ ១៣កើត ខែមិគសិរ ឆ្នាំមមី អដ្ឋស័ក ពុទ្ធសករាជ ២៥៧០\\[2pt]
  ព្រះក្រសាំង, ថ្ងៃទី២២ ខែធ្នូ ឆ្នាំ២០២៦\\[4pt]
  {\muol គ្រូបង្រៀន}\\[32pt]
  {\muol ហុក ប៊ុនថែន}
\end{minipage}
\hspace*{0.8cm}

\end{document}
